\documentclass[11pt]{article}
\usepackage[a4paper,margin=1in]{geometry}
\usepackage{amsmath,amssymb,booktabs,enumitem,mathtools}
\usepackage[T1]{fontenc}
\usepackage[utf8]{inputenc}
\usepackage{lmodern}
\usepackage{hyperref}

\title{Detailed Note on the Tucker Rank-Selection Algorithm\\for ROI10 Tensor Regression}
\author{}
\date{\today}

\begin{document}
\maketitle

\section{Purpose}
This note documents, in detail, the tensor rank-selection algorithm implemented in
\texttt{H:\textbackslash HBN-EEG\textbackslash code\textbackslash run\_roi10\_tucker\_tensor\_regression\_cv.m}.
The goal is to explain:
\begin{itemize}[leftmargin=2em]
\item how the EEG ROI-band summaries are converted into a $10\times 8\times 2$ tensor covariate,
\item how 5-fold cross-validation is performed over the Tucker rank grid,
\item how the fitting routine \texttt{fit\_tucker\_tensor\_regression.m} works internally,
\item how the selected rank is determined, and
\item what the current default script outputs on the real data.
\end{itemize}

\section{Input Data and Target Model}
\subsection{Tensor Construction}
The script reads
\texttt{H:\textbackslash HBN-EEG\textbackslash derivatives\textbackslash bandpower\_8band\_ordered\_preproc\_roi10\textbackslash matlab\textbackslash hbn\_bandpower\_8band\_roi10.mat},
which contains:
\begin{itemize}[leftmargin=2em]
\item \texttt{EO\_roi\_power\_uV2},
\item \texttt{EC\_roi\_power\_uV2},
\item \texttt{age},
\item \texttt{roi\_names},
\item \texttt{band\_names}.
\end{itemize}

For subject $i$, let
\[
X_i^{(EC)} \in \mathbb{R}^{10\times 8}, \qquad
X_i^{(EO)} \in \mathbb{R}^{10\times 8}
\]
denote the raw ROI-by-band power matrices in the EC and EO conditions.

The script first defines a positive floor
\[
\texttt{rawFloor}
=
\max\left(\frac{1}{2}\min\{x:x>0\},\,10^{-12}\right),
\]
where the minimum is taken over all positive entries in the EC and EO arrays.
It then applies the transform
\[
\widetilde X_i^{(EC)}=\log_{10}\!\bigl(\max(X_i^{(EC)},\texttt{rawFloor})\bigr),\qquad
\widetilde X_i^{(EO)}=\log_{10}\!\bigl(\max(X_i^{(EO)},\texttt{rawFloor})\bigr).
\]

The tensor covariate is constructed by stacking the two condition-specific matrices:
\[
\mathcal X_i(:,:,1)=\widetilde X_i^{(EC)},\qquad
\mathcal X_i(:,:,2)=\widetilde X_i^{(EO)}.
\]
Hence each subject contributes a third-order tensor
\[
\mathcal X_i \in \mathbb{R}^{10\times 8\times 2},
\]
with modes corresponding to:
\begin{itemize}[leftmargin=2em]
\item ROI,
\item frequency band,
\item condition (EC / EO).
\end{itemize}

The response is age:
\[
y_i = \texttt{age}_i.
\]
The model fitted by the script is a Tucker low-rank tensor regression model,
\[
y_i \approx \langle \mathcal X_i,\mathcal B\rangle + \varepsilon_i,
\]
where the coefficient tensor $\mathcal B\in\mathbb{R}^{10\times 8\times 2}$ is constrained to have low Tucker rank.

\subsection{Candidate Rank Grid}
The current default rank grid is
\[
r_1\in\{1,\dots,10\},\qquad
r_2\in\{1,\dots,5\},\qquad
r_3\in\{1,2\}.
\]
Therefore the total number of candidate rank triples is
\[
10\times 5\times 2 = 100.
\]
The script constructs the full grid
\[
\mathcal R=\{(r_1,r_2,r_3)\}
\]
using \texttt{ndgrid}.

\section{Outer Cross-Validation Procedure}
\subsection{Fold Assignment}
Let $N$ denote the sample size. The script fixes the random seed
\[
\texttt{split\_seed}=20260321
\]
and generates a random permutation $\pi$ of the subject indices. The fold labels are then assigned according to
\[
\texttt{foldId}(\pi_j)=(j-1)\bmod 5 + 1.
\]
Thus each subject is assigned to one of five folds, and the same fold assignment is reused for all candidate ranks.

\subsection{Evaluation of a Fixed Rank Triple}
Fix a candidate Tucker rank
\[
\mathbf r=(r_1,r_2,r_3).
\]
For each fold $k=1,\dots,5$:
\begin{itemize}[leftmargin=2em]
\item the training set is all observations with \texttt{foldId $\neq k$},
\item the validation set is all observations with \texttt{foldId $= k$}.
\end{itemize}

The script fits the model on the training set via
\[
\texttt{fit\_tucker\_tensor\_regression}(X_{\mathrm{train}},y_{\mathrm{train}},\mathbf r),
\]
and obtains validation predictions via
\[
\texttt{predict\_tucker\_tensor\_regression}(\widehat{\mathcal B},X_{\mathrm{val}}).
\]

The validation loss is
\[
\mathrm{MSE}_k(\mathbf r)
=
\frac{1}{n_k}\sum_{i\in\mathcal V_k}(y_i-\widehat y_i)^2.
\]
The script also reports
\[
R_k^2(\mathbf r)
=
1-
\frac{\sum_{i\in\mathcal V_k}(y_i-\widehat y_i)^2}
{\sum_{i\in\mathcal V_k}(y_i-\bar y_{\mathcal V_k})^2},
\]
where $\bar y_{\mathcal V_k}$ is the mean of the validation responses within fold $k$.

\subsection{Aggregation Across Folds}
For each candidate rank $\mathbf r$, the script computes
\[
\overline{\mathrm{MSE}}(\mathbf r)
=
\frac{1}{5}\sum_{k=1}^5 \mathrm{MSE}_k(\mathbf r),
\qquad
\overline{R^2}(\mathbf r)
=
\frac{1}{5}\sum_{k=1}^5 R_k^2(\mathbf r).
\]
It also stores the empirical standard deviations of fold-wise MSE and fold-wise $R^2$.

\subsection{Selection Rule}
The summary table is sorted by the following priority:
\begin{enumerate}[leftmargin=2em,label=\arabic*.]
\item \texttt{mean\_val\_mse} in ascending order,
\item \texttt{rank\_sum}=$r_1+r_2+r_3$ in ascending order,
\item \texttt{rank1} in ascending order,
\item \texttt{rank2} in ascending order,
\item \texttt{rank3} in ascending order.
\end{enumerate}
Hence the selected rank is the first row of the sorted summary table.
Equivalently:
\begin{quote}
Choose the rank triple with the smallest mean validation MSE; if there is a tie, prefer the lower total rank, and then the lexicographically smaller rank triple.
\end{quote}

\section{Internal Fitting Routine}
\subsection{Centering Within the Training Fold}
Suppose the training set has size $n$. The fitting routine first computes
\[
\bar{\mathcal X}
=
\frac{1}{n}\sum_{i=1}^n \mathcal X_i,
\qquad
\bar y
=
\frac{1}{n}\sum_{i=1}^n y_i,
\]
and centers the training data:
\[
\mathcal X_i^c=\mathcal X_i-\bar{\mathcal X},
\qquad
y_i^c=y_i-\bar y.
\]
The stored model retains \texttt{Xmean} and \texttt{ymean}. At prediction time, new samples are centered using the \emph{training-fold mean}, not the validation-fold mean.

\subsection{Vectorized Representation}
The centered tensor observations are reshaped into a matrix
\[
Z\in\mathbb{R}^{n\times (10\cdot 8\cdot 2)}.
\]
The optimization target becomes
\[
\min_{\mathcal B:\,\operatorname{rank}_{\mathrm{Tucker}}(\mathcal B)\le (r_1,r_2,r_3)}
\frac{1}{2n}\|Z\,\mathrm{vec}(\mathcal B)-y^c\|_2^2.
\]

\subsection{Initialization}
The script first solves the unconstrained least-squares problem
\[
\beta_0
=
\arg\min_{\beta}\|Z\beta-y^c\|_2^2.
\]
Implementation details:
\begin{itemize}[leftmargin=2em]
\item it uses \texttt{lsqminnorm(Z, ycentered)} if available;
\item otherwise it falls back to \texttt{pinv(Z) * ycentered}.
\end{itemize}
The vector $\beta_0$ is reshaped into a tensor $\mathcal B^{(0)}$ and projected to the candidate Tucker rank.

\subsection{Projection to Fixed Tucker Rank}
Given a tensor $\mathcal B$, the script computes the leading left singular vectors of each mode unfolding:
\[
U_1=\operatorname{SVD}_{r_1}(B_{(1)}),\qquad
U_2=\operatorname{SVD}_{r_2}(B_{(2)}),\qquad
U_3=\operatorname{SVD}_{r_3}(B_{(3)}).
\]
Then it forms the core tensor
\[
\mathcal G
=
\mathcal B\times_1 U_1^\top\times_2 U_2^\top\times_3 U_3^\top
\]
and reconstructs the projected tensor
\[
\Pi_{\mathbf r}(\mathcal B)
=
\mathcal G\times_1 U_1\times_2 U_2\times_3 U_3.
\]
This is the projection operator implemented in \texttt{project\_tucker\_tensor}.

\subsection{Projected Gradient Descent}
After initialization, the routine iterates
\[
\mathcal B^{(t+\frac12)}
=
\mathcal B^{(t)}-\eta \nabla f(\mathcal B^{(t)}),
\]
followed by projection,
\[
\mathcal B^{(t+1)}
=
\Pi_{\mathbf r}\!\left(\mathcal B^{(t+\frac12)}\right),
\]
where
\[
f(\mathcal B)
=
\frac{1}{2n}\|Z\,\mathrm{vec}(\mathcal B)-y^c\|_2^2.
\]
The gradient is
\[
\nabla f(\mathcal B)
=
\operatorname{reshape}\!\left(\frac{1}{n}Z^\top(Z\,\mathrm{vec}(\mathcal B)-y^c)\right).
\]

\subsection{Step Size}
Let $s_{\max}(Z)$ denote the largest singular value of the design matrix $Z$. The script defines
\[
L
=
\max\left(\frac{s_{\max}(Z)^2}{n},10^{-8}\right),
\qquad
\eta=\frac{1}{L}.
\]
Thus the step size is the inverse of a Lipschitz upper bound for the quadratic loss gradient.

\subsection{Stopping Rule}
At iteration $t$, the algorithm monitors:
\begin{align*}
\text{relative coefficient change}
&=
\frac{\|\mathcal B^{(t+1)}-\mathcal B^{(t)}\|_F}
{\max(\|\mathcal B^{(t)}\|_F,10^{-12})},\\[4pt]
\text{relative objective change}
&=
\frac{|f(\mathcal B^{(t+1)})-f(\mathcal B^{(t)})|}
{\max(|f(\mathcal B^{(t)})|,1)}.
\end{align*}
The routine stops when either quantity is below \texttt{Tol}.
The current rank-selection script uses
\[
\texttt{MaxIter}=4000,\qquad \texttt{Tol}=10^{-6}.
\]

\subsection{Prediction Formula}
Let $\widehat{\mathcal B}$ denote the fitted coefficient tensor, and let $\bar{\mathcal X}$ and $\bar y$ be the training-fold means stored in the model. For a new tensor covariate $\mathcal X$, the prediction is
\[
\widehat y
=
\bar y + \langle \mathcal X-\bar{\mathcal X},\widehat{\mathcal B}\rangle.
\]
This ensures that training and validation use the same centering convention.

\section{Default Output on the Real Data}
The default output directory is
\texttt{H:\textbackslash HBN-EEG\textbackslash derivatives\textbackslash bandpower\_8band\_ordered\_preproc\_roi10\textbackslash tucker\_tensor\_regression\_age\_log10\_ec\_eo\_cv}.

The top 10 rows of the current \texttt{tucker\_tensor\_regression\_cv\_summary.csv} are:
\begin{center}
\begin{tabular}{ccccccc}
\toprule
$r_1$ & $r_2$ & $r_3$ & $r_1+r_2+r_3$ & mean val MSE & std val MSE & mean val $R^2$ \\
\midrule
6  & 2 & 2 & 10 & 7.9627 & 2.1841 & 0.2337 \\
5  & 2 & 2 & 9  & 7.9708 & 2.1749 & 0.2330 \\
3  & 2 & 2 & 7  & 8.0475 & 1.9957 & 0.2260 \\
7  & 2 & 2 & 11 & 8.0672 & 2.1691 & 0.2236 \\
4  & 2 & 2 & 8  & 8.0689 & 2.1515 & 0.2236 \\
9  & 2 & 2 & 13 & 8.0698 & 2.1733 & 0.2237 \\
8  & 2 & 2 & 12 & 8.0752 & 2.1699 & 0.2230 \\
10 & 2 & 2 & 14 & 8.0790 & 2.1696 & 0.2227 \\
3  & 2 & 1 & 6  & 8.2685 & 2.1458 & 0.2058 \\
4  & 2 & 1 & 7  & 8.2881 & 2.1243 & 0.2037 \\
\bottomrule
\end{tabular}
\end{center}

Under the current default implementation, the selected rank is therefore
\[
\widehat{\mathbf r}_{CV}=(6,2,2).
\]
This is the rank selected by the script itself. A later inference analysis may override this and use $(4,2,2)$ for stability considerations, but that is a separate modeling choice and not part of the cross-validation routine described here.

\section{Files Produced by the Script}
The rank-selection script writes the following main outputs:
\begin{itemize}[leftmargin=2em]
\item \texttt{tucker\_tensor\_regression\_cv\_details.csv}: fold-by-fold validation results for every candidate rank;
\item \texttt{tucker\_tensor\_regression\_cv\_summary.csv}: rank-level summary table after aggregating over folds;
\item \texttt{tucker\_tensor\_regression\_best\_model.mat}: final model refit on the full sample using the selected rank;
\item \texttt{tucker\_tensor\_regression\_best\_model\_coefficients.csv}: elementwise export of the final coefficient tensor;
\item \texttt{tucker\_tensor\_regression\_cv\_mean\_mse.png}: CV heatmap over the rank grid;
\item \texttt{tucker\_tensor\_regression\_best\_model\_coefficients.png}: image of the final coefficient slices;
\item \texttt{tucker\_tensor\_regression\_summary.txt}: text summary of the selected rank and final fit.
\end{itemize}

\section{Algorithm in Compact Form}
The current implementation can be summarized as follows:
\begin{enumerate}[leftmargin=2em,label=\arabic*.]
\item Read the ROI10 EEG data and construct $\mathcal X_i\in\mathbb{R}^{10\times 8\times 2}$ and $y_i$.
\item Generate a fixed 5-fold split using the preset random seed.
\item For each candidate Tucker rank $\mathbf r\in\mathcal R$:
\begin{enumerate}[leftmargin=2em,label*=\arabic*.]
\item For each fold:
\begin{enumerate}[leftmargin=2em,label*=\arabic*.]
\item Center the training tensors and training responses.
\item Obtain an unconstrained least-squares initializer.
\item Project the initializer to Tucker rank $\mathbf r$.
\item Iterate projected gradient descent until convergence.
\item Predict the validation fold using the training-fold centering constants.
\item Record validation MSE and validation $R^2$.
\end{enumerate}
\item Aggregate the 5 validation losses for this rank.
\end{enumerate}
\item Sort all candidate ranks by mean validation MSE, then by lower total rank.
\item Select the first rank in that sorted table.
\item Refit the model on the full sample using the selected rank and export the results.
\end{enumerate}

\section{Recommended One-Sentence Description}
If a concise description is needed in a paper or report, the algorithm can be summarized as:
\begin{quote}
We form a $10\times 8\times 2$ tensor covariate from the log-transformed EC and EO ROI-band power matrices, fit Tucker low-rank tensor regression models over a prespecified rank grid, and select the Tucker rank by 5-fold cross-validation using the mean validation MSE, with ties broken in favor of lower total rank.
\end{quote}

\end{document}
